% ============================================================
% setup/format.tex —— 正文排版规范（按学位论文要求集中配置）
%
% 本文件由 main.tex 通过 \input 引入（放在 commands 之后）。
% 覆盖内容：
%   1) 章节/节/小节标题：字号、字体、对齐、上下间距、序号与题名间距
%   2) 正文段落：宋体小四、英文 Times New Roman、两端对齐、
%                首行缩进 2 汉字符、固定行距 20 磅、段前段后 0
%   3) 图/表序号与题名：宋体五号、居中、单倍行距、序号与题名空 1 汉字符
%   4) 表达式编号：右对齐、(n) 形式、Times New Roman 五号
%
% 字号对照：三号≈16pt / 四号≈14pt / 小四≈12pt / 五号≈10.5pt
% ============================================================

% ---------- 1. 标题格式（ctex 标题机制）----------
\ctexset{
   % 各章标题：黑体三号(16pt)（黑体本身即粗，不再叠加 \bfseries 以免伪粗警告）居中，单倍行距，上下空 2 行
  % 章序号用中文数字（第一章），序号与题名之间空一个汉字符
  chapter = {
    name      = {第,章},
    number    = \chinese{chapter},   % 若学校要求"第1章"阿拉伯数字，改为 \arabic{chapter}
    format    = \centering\heiti\zihao{3},
    aftername = \quad,               % 序号与题名间空一个汉字符
    beforeskip = 2\baselineskip,     % 上空约 2 行
    afterskip  = 2\baselineskip,     % 下空约 2 行
    fixskip   = true,                % 固定间距（无弹性）
  },
  % 一级节标题：黑体四号(14pt)（黑体不加 \bfseries）居中，单倍行距，上下空 1 行
  section = {
    format    = \centering\heiti\zihao{4},
    aftername = \quad,
    beforeskip = \baselineskip,      % 上空约 1 行
    afterskip  = \baselineskip,      % 下空约 1 行
    fixskip   = true,
  },
  % 二级节标题：黑体小四(12pt)居左，单倍行距，上下空 1 行
  subsection = {
    format    = \raggedright\heiti\zihao{-4},
    aftername = \quad,
    beforeskip = \baselineskip,
    afterskip  = \baselineskip,
    fixskip   = true,
  },
  % 三级节标题（如使用）：规范表3.8未单独规定，按 CY/T 35-2001 顺延为
  % 黑体小四、顶格（序号与题名间不额外加空格），上下空约半行
  subsubsection = {
    format    = \raggedright\heiti\zihao{-4},
    aftername = {},
    beforeskip = 0.5\baselineskip,
    afterskip  = 0.5\baselineskip,
    fixskip   = true,
  },
}

% ---------- 2. 正文段落 ----------
% 宋体小四由文档类 zihao=-4 统一控制；英文 Times New Roman 由 fonts.tex 控制
% 其余：两端对齐、首行缩进 2 汉字符、固定行距 20 磅、段前段后 0
\setlength{\parindent}{2\ccwd}      % 首行左缩进 2 个汉字符
\setlength{\parskip}{0pt}           % 段前段后 0（段与段之间不额外留空）
% 固定值行距 20 磅：直接控制 \baselineskip。
% 注意 1：setspace 并没有 \exactlineskipmode 命令，也没有 [exact] 选项
%         （会报 "Unknown option `exact'"），不可依赖它。
% 注意 2：直接改 \baselineskip 会被字号命令（\zihao{5} 等）重置。
% 方案：用 ctex 官方钩子 \ctex_patch_cmd 把"固定 20pt"追加进各字号的
%       内部定义——每次切换字号都会重新执行它们，行距恒定。
%       \lineskiplimit=-1000pt 是 TeX"固定行距"模式，等效 Word"固定值"。
\newcommand{\fixLineskip}[1]{%
  \setlength{\baselineskip}{#1}%
  \setlength{\lineskip}{\z@}%
  \setlength{\lineskiplimit}{-1000\p@}%  % TeX 固定行距模式
}
\makeatletter
% \normalsize/\small/\footnotesize 分别调用内部 \@normalsize/\@small/\@footnotesize。
% LaTeX 内核的 \addto 要求第二参数是"已展开的命令名"，因此先把
% \fixLineskip{20pt} 展开到 \next（用 \dimexpr...\relax 隔离 \p@ 的 @），再追加。
\expandafter\def\expandafter\next\expandafter{\expandafter\fixLineskip\expandafter{\dimexpr20\p@\relax}}
\addto\@normalsize{\next}
\@ifundefined{@small}{}{\addto\@small{\next}}
\@ifundefined{@footnotesize}{}{\addto\@footnotesize{\next}}
\makeatother
\AtBeginDocument{\normalsize}   % 使上述设置立即生效

% ---------- 3. 图 / 表序号与题名 ----------
% 宋体五号(10.5pt)居中，单倍行距，序号与题名之间空一个汉字符
% 图序置于图下方（默认），表序置于表上方
\usepackage{caption}
% 注意：caption 的 font 键不接受 \songti、\zihao{5} 这类"声明式"命令，
% 必须使用 caption 提供的受控字体/字号名（标准尺寸名或 \DeclareCaptionFont
% 自定义注册名）。中文规范"题注宋体五号居中"实现如下：
%   wuhao  = 严格五号（10.5pt/13pt 单倍行距）；
%   zhsongcn = ctex 定义的宋体族控制序列 \songti。caption 的 font 键
%            按名字展开字体命令，因此用 \csname songti\endcsname 间接
%            引用；songti 由 ctex（scheme=chinese）定义。
\DeclareCaptionFont{wuhao}{\fontsize{10.5pt}{13pt}\selectfont}
% 宋体族名：统一使用 ctex 定义的 CJKfamily "zhsong"
% ctex 定义的 CJKfamily "zhsong"（windows/mac/linux/fandol 均注册此名）。
% caption 在正常文档阶段展开，此时 \songti 已由 ctex 定义。
\DeclareCaptionFont{zhsongcn}{\csname songti\endcsname}  % 宋体（ctex 提供）
\captionsetup{
  font         = {zhsongcn,wuhao},      % 宋体五号
  labelfont    = {zhsongcn,wuhao},      % 图/表序号同样宋体五号
  justification = centering,          % 居中
  labelsep     = quad,                % 序号与题名之间空一个汉字符
  skip         = 6pt,                 % 图/表与题名的间距
}
\captionsetup[table]{position=above}   % 表序在表上方
\captionsetup[figure]{position=below}  % 图序在图下方

% ---------- 4. 表达式编号 ----------
% 右对齐、(n) 形式、Times New Roman 五号（默认已带圆括号，仅改字体与字号）
\makeatletter
\def\maketag@@@#1{\hbox{\m@th\normalfont\zihao{5}#1}}
\makeatother

% ---------- 5. 目录排版 ----------
% 标题"目录"：沿用 ctexset 的 chapter 样式（黑体三号加粗居中，单倍行距），无需额外设置。
% 目录条目要求：
%   各章目录  ：宋体四号(14pt)，固定行距 20 磅，两端对齐，页码右对齐，左缩进 0
%   一级节    ：宋体小四(12pt)，固定行距 20 磅，两端对齐，页码右对齐，左缩进 1 汉字符
%   二级节    ：宋体小四(12pt)，固定行距 20 磅，两端对齐，页码右对齐，左缩进 2 汉字符
\usepackage{titletoc}
\titlecontents{chapter}[0pt]{%                        % 左缩进 0
    \zihao{4}\songti\baselineskip=20pt%               % 宋体四号 + 固定 20 磅行距
  }{%
    \thecontentslabel\quad%                           % 编号（第一章）+ 空一个汉字符
  }{%                                                  % 无编号条目（如参考文献/致谢）
    {}%
  }{%
    \hfill\contentspage%                               % 页码右对齐（无点线）
  }
\titlecontents{section}[1\ccwd]{%                      % 左缩进 1 汉字符
    \zihao{-4}\songti\baselineskip=20pt%              % 宋体小四 + 固定 20 磅行距
  }{%
    \thecontentslabel\quad%
  }{%
    {}%
  }{%
    \hfill\contentspage%
  }
\titlecontents{subsection}[2\ccwd]{%                   % 左缩进 2 汉字符
    \zihao{-4}\songti\baselineskip=20pt%
  }{%
    \thecontentslabel\quad%
  }{%
    {}%
  }{%
    \hfill\contentspage%
  }

% 图目录 / 表目录：标题沿用 chapter 样式（黑体三号加粗居中）；
% 条目按一级节样式：宋体小四(12pt)、固定 20 磅、页码右对齐、左缩进 1 汉字符
\titlecontents{figure}[1\ccwd]{%
    \zihao{-4}\songti\baselineskip=20pt%
  }{%
    \thecontentslabel\quad%
  }{}{%
    \hfill\contentspage%
  }
\titlecontents{table}[1\ccwd]{%
    \zihao{-4}\songti\baselineskip=20pt%
  }{%
    \thecontentslabel\quad%
  }{}{%
    \hfill\contentspage%
  }

% ---------- 6. 中英文摘要 ----------
% 标题：中文"摘要"黑体小二(18pt)加粗居中；英文"Abstract" Times New Roman 小二(18pt)加粗居中；单倍行距
% 正文：中文宋体小四 / 英文 TNR 小四，固定 20 磅，段前段后 0（全局已设）
% 关键词："关键词"三字加粗（黑体）/ "Key Words"两词加粗（TNR），字号同正文
\newcommand{\cnabstracttitle}{%
  \begin{center}\heiti\zihao{-2} 摘要\end{center}%
  \vspace{0.5\baselineskip}\zihao{-4}\songti\noindent\ignorespaces
}
\newcommand{\enabstracttitle}{%
  \begin{center}\bfseries\zihao{-2} Abstract\end{center}%
  \vspace{0.5\baselineskip}\normalfont\zihao{-4}\noindent\ignorespaces
}

% ---------- 7. 纸张与页眉页脚（中国地质大学硕士论文格式）----------
% 纸张：A4（210×297），白色
% 页边距：上/下 3cm，左/右 3cm，装订线 0cm
% 页眉距顶 2.5cm（headheight=2.5cm，headsep=0.5cm）；页脚距底 2.0cm（footskip=2.0cm）
\geometry{
  a4paper,
  top        = 3cm,
  bottom     = 3cm,
  left       = 3cm,
  right      = 3cm,
  headheight = 2.5cm,   % 页眉区域高度（页眉底边距顶 2.5cm，贴顶）
  headsep    = 0.5cm,   % 页眉底边到正文的距离
  % 注意：geometry 没有 footheight 键（正确键名为 bottomheight / footskip），
  % 原 footheight=1.0cm 属无效选项会被忽略并警告，此处删除；
  % 页脚位置由 footskip 控制。
  footskip   = 2.0cm,   % 正文底边到页脚基线的距离（页脚距底约 2.0cm，贴底）
  bindingoffset = 0cm   % 装订线 0cm
}

% 页眉页脚样式：正文（第一章起）使用 running header
%   奇数页居中：中国地质大学<学位类别>
%   偶数页居中：作者姓名：论文题目
%   页码外侧（奇→右、偶→左），宋体五号(10.5pt)
%   页眉之下有下划线（headrule）
\fancyhf{}                                     % 清空默认页眉页脚
\renewcommand{\headrulewidth}{0.4pt}           % 页眉下划线
\renewcommand{\footrulewidth}{0pt}             % 页脚无线
\fancyhead[CO]{\songti\zihao{5} 中国地质大学\thesisCategory}      % 奇数页：学位类别
\fancyhead[CE]{\songti\zihao{5} \thesisAuthor：\thesisTitle}      % 偶数页：作者：题目
\fancyfoot[RO]{\songti\zihao{5} \thepage}      % 奇数页页码在右侧（外侧）
\fancyfoot[LE]{\songti\zihao{5} \thepage}      % 偶数页页码在左侧（外侧）

% 英文摘要页眉：按规范"Abstract 部分用 Times New Roman 五号"，标题居中（TNR）
\fancypagestyle{abstracten}{%
  \fancyhf{}%
  \renewcommand{\headrulewidth}{0.4pt}%
  \fancyhead[C]{\zihao{5} \thesisTitleEN}%      % 英文标题（TNR，主字体即 Times New Roman）
  \fancyfoot[C]{\songti\zihao{5} \thepage}%     % 页码仍按规范用宋体五号
}
